\documentclass[answers,12pt,addpoints]{exam}
\usepackage{import}
 
\import{C:/Users/prana/OneDrive/Desktop/MathNotes}{style.tex}
 
% Packages used specifically in this homework file.
% (Reloading amsmath/amssymb is harmless if style.tex already loads them.)
\usepackage{amsmath,amssymb}
\usepackage{booktabs}
\usepackage{listings}
\usepackage{xcolor}
 
% Header
\newcommand{\name}{Pranav Tikkawar}
\newcommand{\course}{Stats 596}
\newcommand{\assignment}{Homework 1}
\author{\name}
\title{\course\ -- \assignment}
\date{Fall 2026}
 
% Fallbacks: these take effect only if style.tex does not already define them.
\providecommand{\R}{\mathbb{R}}
\providecommand{\E}{\operatorname{E}}
\providecommand{\argmin}{\operatorname*{arg\,min}}
 
% Local notation
\newcommand{\trans}{^{\mathsf T}}
\newcommand{\normtwo}[1]{\left\lVert #1\right\rVert_2}
\newcommand{\col}{\mathcal{C}}
\newcommand{\nul}{\mathcal{N}}
\newcommand{\ind}{\mathbf{1}}   % indicator function
\newcommand{\ones}{\mathbf{1}}  % vector of ones
 
\definecolor{codegreen}{rgb}{0,0.45,0.15}
\definecolor{codepurple}{rgb}{0.50,0,0.50}
\definecolor{codegray}{rgb}{0.40,0.40,0.40}
\definecolor{codeback}{rgb}{0.97,0.97,0.97}
\lstdefinestyle{Rstyle}{
  language=R,
  backgroundcolor=\color{codeback},
  basicstyle=\ttfamily\small,
  keywordstyle=\color{blue},
  commentstyle=\color{codegreen},
  stringstyle=\color{codepurple},
  numberstyle=\tiny\color{codegray},
  numbers=left,
  numbersep=6pt,
  breaklines=true,
  columns=fullflexible,
  showstringspaces=false,
  keepspaces=true,
  tabsize=2
}
\lstset{style=Rstyle}
 
\begin{document}
\maketitle
 
\begin{questions}
 
%==============================================================================
\question \textbf{Exercise 1.3.}
Let $A,B,C,$ and $D$ be matrices of appropriate dimensions and write
\[
M=\begin{pmatrix}A&B\\C&D\end{pmatrix},
\qquad
\widetilde A=A-BD^{-1}C.
\]
 
\begin{parts}
\part Show that if $D$ and $\widetilde A$ are nonsingular, then $M$ is
nonsingular and
\[
M^{-1}=
\begin{pmatrix}
\widetilde A^{-1} & -\widetilde A^{-1}BD^{-1}\\
-D^{-1}C\widetilde A^{-1} &
D^{-1}+D^{-1}C\widetilde A^{-1}BD^{-1}
\end{pmatrix}.
\]
Give a direct proof using the orthogonalization idea or a variation.
 
\part Show that if $D$ and $M$ are nonsingular, then $\widetilde A$ is
nonsingular and the preceding inverse formula holds.
 
\part Suppose $A,D,$ and $\widetilde A$ are nonsingular.  Show that
\[
\widetilde D=D-CA^{-1}B
\]
is nonsingular and that
\[
\widetilde D^{-1}
=D^{-1}+D^{-1}C\widetilde A^{-1}BD^{-1}.
\]
 
\part As an application, show that if $u\trans u\neq 1$, then $I-uu\trans$
is nonsingular and
\[
(I-uu\trans)^{-1}=I+(1-u\trans u)^{-1}uu\trans.
\]
\end{parts}
 
\begin{solution}
Take $A$ to be $m\times m$ and $D$ to be $q\times q$. The matrix
$\widetilde A=A-BD^{-1}C$ is the Schur complement of $D$ in $M$.
 
\textbf{(a)} Subtracting $BD^{-1}$ times the second block row from the first
removes $B$; this residualizes $A$ on the $D$-block, which is the
orthogonalization idea. In matrix form,
\[
M=LT,\qquad
L=\begin{pmatrix}I&BD^{-1}\\0&I\end{pmatrix},\qquad
T=\begin{pmatrix}\widetilde A&0\\C&D\end{pmatrix},
\]
as can be checked by multiplying out. Both factors are invertible, with
\[
L^{-1}=\begin{pmatrix}I&-BD^{-1}\\0&I\end{pmatrix},\qquad
T^{-1}=\begin{pmatrix}\widetilde A^{-1}&0\\-D^{-1}C\widetilde A^{-1}&D^{-1}\end{pmatrix},
\]
again by direct multiplication. So $M$ is invertible, and multiplying out
$M^{-1}=T^{-1}L^{-1}$ gives the stated formula.
 
\textbf{(b)} $L$ is invertible whenever $D$ is, so $T=L^{-1}M$ is invertible.
If $\widetilde Av=0$ for some $v\neq0$, put $w=-D^{-1}Cv$; then
\[
T\begin{pmatrix}v\\w\end{pmatrix}
=\begin{pmatrix}\widetilde Av\\Cv+Dw\end{pmatrix}=0,
\]
contradicting invertibility of $T$. So $\widetilde A$ is invertible and part
(a) applies.
 
\textbf{(c)} Let $\widetilde D=D-CA^{-1}B$. Eliminating with $A$ as the pivot
instead gives
\[
M=L_AU_A,\qquad
L_A=\begin{pmatrix}I&0\\CA^{-1}&I\end{pmatrix},\qquad
U_A=\begin{pmatrix}A&B\\0&\widetilde D\end{pmatrix}.
\]
Since $D$ and $\widetilde A$ are nonsingular, which are exactly the hypotheses of (a), $M$ is nonsingular. Therefore $U_A=L_A^{-1}M$ is nonsingular too. A block upper-triangular
matrix whose $(1,1)$ block $A$ is invertible can only be invertible if its
$(2,2)$ block is: if $\widetilde Dv=0$ with $v\neq0$, then $w=-A^{-1}Bv$ gives
$U_A\begin{pmatrix}w\\v\end{pmatrix}=0$. Hence $\widetilde D$ is invertible,
and
\[
M^{-1}=U_A^{-1}L_A^{-1}
=\begin{pmatrix}
A^{-1}+A^{-1}B\widetilde D^{-1}CA^{-1}&-A^{-1}B\widetilde D^{-1}\\
-\widetilde D^{-1}CA^{-1}&\widetilde D^{-1}
\end{pmatrix}.
\]
Since $M$ has only one inverse, this $(2,2)$ block equals the one from (a):
$\widetilde D^{-1}=D^{-1}+D^{-1}C\widetilde A^{-1}BD^{-1}$.
 
\textbf{(d)} Apply (c) with $A=1$, $B=u\trans$, $C=u$, and $D=I$. Then
$\widetilde A=1-u\trans u\neq0$ and $\widetilde D=I-uu\trans$, so $I-uu\trans$
is invertible and
\[
(I-uu\trans)^{-1}=I+u(1-u\trans u)^{-1}u\trans=I+(1-u\trans u)^{-1}uu\trans.
\]
\end{solution}
 
%==============================================================================
\question \textbf{Exercise 1.5.}
Let $X$ be an $n\times p$ covariate matrix and let $Y$ be an $n\times1$
response vector.
 
\begin{parts}
\part Show that if $n\geq p$ and $X$ has rank $p$, then $X\trans X$ is
nonsingular and
\[
\widehat\beta=(X\trans X)^{-1}X\trans Y
\]
is the unique solution of
$\min_{\beta\in\R^p}\normtwo{Y-X\beta}^2$.
 
\part Show that if $n\leq p$ and $X$ has rank $n$, then $XX\trans$ is
nonsingular and
\[
\widetilde\beta=X\trans(XX\trans)^{-1}Y
\]
is the unique solution of
\[
\min_{\beta\in\R^p}\normtwo{\beta}
\quad\text{subject to}\quad Y=X\beta.
\]
 
\part Show that, regardless of the rank of $X$,
\[
(X\trans X)^+X\trans Y=X\trans(XX\trans)^+Y.
\]
\end{parts}
 
\begin{solution}
\textbf{(a)} Full column rank means $Xb\neq0$ whenever $b\neq0$, so
$b\trans X\trans Xb=\normtwo{Xb}^2>0$; thus $X\trans X$ is positive definite and
invertible. Let $\widehat r=Y-X\widehat\beta$. Then
$X\trans\widehat r=X\trans Y-X\trans X(X\trans X)^{-1}X\trans Y=0$, so for any $\beta=\widehat\beta+h$,
\[
\normtwo{Y-X\beta}^2=\normtwo{\widehat r-Xh}^2
=\normtwo{\widehat r}^2+\normtwo{Xh}^2,
\]
because the cross term $2h\trans X\trans\widehat r$ vanishes. This is larger
than $\normtwo{\widehat r}^2$ unless $Xh=0$, that is, unless $h=0$.
 
\textbf{(b)} $\operatorname{rank}(X)=n$ means $\operatorname{rank}(X\trans)=n$,
so $\nul(X\trans)=\{0\}$ and $X\trans a\neq0$ for $a\neq0$. Thus
$a\trans XX\trans a=\normtwo{X\trans a}^2>0$ and $XX\trans$ is invertible.
Write $\widetilde\beta=X\trans\widetilde a$ with $\widetilde a=(XX\trans)^{-1}Y$;
then $X\widetilde\beta=Y$. Any other solution of $X\beta=Y$ has the form
$\beta=\widetilde\beta+h$ with $Xh=0$, and
$h\trans\widetilde\beta=(Xh)\trans\widetilde a=0$. Hence
$\normtwo{\beta}^2=\normtwo{\widetilde\beta}^2+\normtwo{h}^2$, which is smallest
only when $h=0$.
 
\textbf{(c)} Take the SVD $X=U\Sigma V\trans$, where $U$ and $V$ are
orthogonal and $\Sigma$ is $n\times p$ and diagonal, with positive singular
values $s_1,\ldots,s_r$ followed by zeros. Since $U$ and $V$ are orthogonal,
$(X\trans X)^+=V(\Sigma\trans\Sigma)^+V\trans$ and
$(XX\trans)^+=U(\Sigma\Sigma\trans)^+U\trans$, so
\[
(X\trans X)^+X\trans=V(\Sigma\trans\Sigma)^+\Sigma\trans U\trans,
\qquad
X\trans(XX\trans)^+=V\Sigma\trans(\Sigma\Sigma\trans)^+U\trans.
\]
Every matrix between $V$ and $U\trans$ is diagonal, so it is enough to compare
diagonal entries. For $s_j>0$ they are $s_j^{-2}s_j=s_j^{-1}$ and
$s_js_j^{-2}=s_j^{-1}$; for $s_j=0$ both are $0$. So the two sides are equal
(both are $X^+$). With full column rank this is the estimator from (a), and
with full row rank it is the minimum-norm solution from (b).
\end{solution}
 
%==============================================================================
\question \textbf{Exercise 1.6.}
Consider the linear regression model
\[
Y_i\sim X_i\trans\beta,
\qquad
X_i=(1,X_{i,1:p}\trans)\trans,
\qquad i=1,\ldots,n.
\]
Let $X$ be the $n\times(p+1)$ design matrix, and suppose
$\widehat\beta=(X\trans X)^{-1}X\trans Y$.  Show that
$\widehat\beta=(\widehat\beta_0,\widehat\beta_{1:p}\trans)\trans$, where
\[
\widehat\beta_0
=\overline Y-\overline X_{1:p}\trans\widehat\beta_{1:p}
\]
and
\[
\widehat\beta_{1:p}
=(\widetilde X_{1:p}\trans\widetilde X_{1:p})^{-1}
 \widetilde X_{1:p}\trans Y,
\]
with $\widetilde X_{1:p}$ denoting the matrix of centered covariates.
 
\begin{solution}
Write $X=(\ones,X_{1:p})$, and assume $X$ has full column rank
(needed for $(X\trans X)^{-1}$ to exist). Let $C=I-\frac1n\ones\ones\trans$ be
the centering matrix, so $CY=Y-\ones\overline Y$ and
$\widetilde X_{1:p}=CX_{1:p}=X_{1:p}-\ones\overline X_{1:p}\trans$. Note that $C$ is symmetric
and $C^2=C$.
 
\textbf{Intercept.} For fixed $b$, $\normtwo{Y-\ones\beta_0-X_{1:p}b}^2$ is a
quadratic in $\beta_0$, and setting its derivative
$-2\ones\trans(Y-\ones\beta_0-X_{1:p}b)$ to zero gives
\[
\widehat\beta_0(b)=\overline Y-\overline X_{1:p}\trans b .
\]
Since $\widehat\beta$ minimizes over $(\beta_0,b)$ jointly, its intercept must
be the best intercept for its own slope, so
\[
\widehat\beta_0=\overline Y-\overline X_{1:p}\trans\widehat\beta_{1:p}.
\]
 
\textbf{Slope.} Plugging in $\widehat\beta_0(b)$, the residual becomes
\[
Y-\ones\widehat\beta_0(b)-X_{1:p}b
=(Y-\ones\overline Y)-(X_{1:p}-\ones\overline X_{1:p}\trans)b
=CY-\widetilde X_{1:p}b .
\]
So $\widehat\beta_{1:p}$ minimizes $\normtwo{CY-\widetilde X_{1:p}b}^2$, which is a
no-intercept regression of $CY$ on $\widetilde X_{1:p}$.
 
This regression has a unique solution because $\widetilde X_{1:p}$ has full column
rank: if $\widetilde X_{1:p}b=0$, then $X_{1:p}b=\ones(\overline X_{1:p}\trans b)$ is
constant, so $(\ones,X_{1:p})\,(-\overline X_{1:p}\trans b,\ b\trans)\trans=0$, and full
rank of $X$ forces $b=0$. By Exercise 1.5(a),
\[
\widehat\beta_{1:p}
=(\widetilde X_{1:p}\trans\widetilde X_{1:p})^{-1}\widetilde X_{1:p}\trans CY
=(\widetilde X_{1:p}\trans\widetilde X_{1:p})^{-1}\widetilde X_{1:p}\trans Y,
\]
where the last step uses
$\widetilde X_{1:p}\trans C=X_{1:p}\trans C^2=\widetilde X_{1:p}\trans$.
\end{solution}
 
%==============================================================================
\question \textbf{Exercise 1.7.}
Consider
\[
Y_i\sim X_i\trans\beta=Z_i\theta+W_i\trans\alpha,
\qquad X=(Z,W),
\qquad
\widehat\beta=(\widehat\theta,\widehat\alpha\trans)\trans.
\]
Let $R_Y$ be the residual vector from regressing $Y$ on $W$, and let $R_Z$
be the residual vector from regressing $Z$ on $W$.  Define
\[
\widehat\sigma^2
=\frac{\normtwo{Y-X\widehat\beta}^2}{n-p},
\qquad
\widetilde\sigma^2
=\frac{\normtwo{R_Y-R_Z\widehat\theta}^2}{n-1}.
\]
 
\begin{parts}
\part Show that
$\widetilde\sigma^2=\frac{n-p}{n-1}\widehat\sigma^2$.
 
\part Show that $(R_Z\trans R_Z)^{-1}$ equals the $(1,1)$ element of
$(X\trans X)^{-1}$.
 
\part Let
$D=\operatorname{diag}(\widehat r_1^2,\ldots,\widehat r_n^2)$, where
$\widehat r_i=Y_i-X_i\trans\widehat\beta$.  Show that
\[
(R_Z\trans R_Z)^{-1}(R_Z\trans D R_Z)(R_Z\trans R_Z)^{-1}
\]
equals the $(1,1)$ element of
\[
(X\trans X)^{-1}(X\trans D X)(X\trans X)^{-1}.
\]
\end{parts}
 
\begin{solution}
Here $Z$ is $n\times1$, $W$ is $n\times(p-1)$, and $X=(Z,W)$ has full column
rank $p$. Let $M_W=I-W(W\trans W)^{-1}W\trans$, which is symmetric with
$M_W^2=M_W$; then $R_Y=M_WY$ and $R_Z=M_WZ$. Note $R_Z\neq0$, since otherwise
$Z\in\col(W)$ and $X$ would not have full rank.
 
\textbf{(a)} For fixed $\theta$, the best $\alpha$ comes from regressing
$Y-Z\theta$ on $W$, namely
$\widehat\alpha(\theta)=(W\trans W)^{-1}W\trans(Y-Z\theta)$, and the residual is
\[
Y-Z\theta-W\widehat\alpha(\theta)=M_W(Y-Z\theta)=R_Y-R_Z\theta .
\]
The joint minimizer must use the best $\alpha$ for its own $\theta$ (otherwise
replacing $\widehat\alpha$ by $\widehat\alpha(\widehat\theta)$ would strictly
lower the criterion), so $\widehat\alpha=\widehat\alpha(\widehat\theta)$, and
$\widehat\theta$ minimizes $\normtwo{R_Y-R_Z\theta}^2$. Hence the two
regressions have identical residuals,
\[
Y-X\widehat\beta=R_Y-R_Z\widehat\theta ,
\]
which is the Frisch--Waugh--Lovell theorem. So both variance estimates have the
same numerator $(n-p)\widehat\sigma^2$, and
$\widetilde\sigma^2=\frac{n-p}{n-1}\widehat\sigma^2$.
 
\textbf{(b)} Partition
\[
X\trans X=\begin{pmatrix}Z\trans Z&Z\trans W\\ W\trans Z&W\trans W\end{pmatrix}.
\]
Both $X\trans X$ and $W\trans W$ are nonsingular, so by Exercise 1.3 the
$(1,1)$ entry of $(X\trans X)^{-1}$ is the inverse of the Schur complement
\[
Z\trans Z-Z\trans W(W\trans W)^{-1}W\trans Z=Z\trans M_WZ=R_Z\trans R_Z .
\]
 
\textbf{(c)} Let $s=R_Z\trans R_Z$. By the block inverse formula in Exercise
1.3, the first column of $(X\trans X)^{-1}$ is
\[
q=\begin{pmatrix}s^{-1}\\ -(W\trans W)^{-1}W\trans Z\,s^{-1}\end{pmatrix},
\]
so
\[
Xq=s^{-1}\bigl(Z-W(W\trans W)^{-1}W\trans Z\bigr)=s^{-1}R_Z .
\]
Since $(X\trans X)^{-1}$ is symmetric, the $(1,1)$ entry of the sandwich is
\[
q\trans X\trans DXq=(Xq)\trans D(Xq)
=(R_Z\trans R_Z)^{-1}(R_Z\trans DR_Z)(R_Z\trans R_Z)^{-1}.
\]
By (a) the two regressions have the same residuals, so $D$ is the same matrix
in both.
\end{solution}
 
%==============================================================================
\question \textbf{Exercise 1.8.}
Generate the data using the following R code:
\begin{lstlisting}
library(MASS)
set.seed(0)
 
n <- 200
p <- 4
mu <- rep(0,p)
msig <- toeplitz(.8^(1:p-1))
sig <- .5
 
x <- mvrnorm(n, mu, msig)
y <- x[,1] + pmax(0, apply(x[,2:p],1,sum))^2 / 3 +
     rnorm(n, 0, sig)
cbind(y,x)[1:3,]
\end{lstlisting}
 
\begin{parts}
\part Report the first three rows.  Fit a linear regression of $y$ on $x$
with an intercept, and report all coefficients, degrees-of-freedom-adjusted
model-based standard errors, and model-robust standard errors.
 
\part Compute the same quantities using residual-on-residual regressions.
 
\part Find the population target values of all coefficients analytically and
numerically.
\end{parts}
 
\begin{solution}
The covariance matrix is $\Sigma_{jk}=0.8^{|j-k|}$: in R, \texttt{1:p-1} means
\texttt{(1:p)-1}, which is \texttt{0:3}.
 
\textbf{(a)} The first three generated rows are
\begin{center}
\begin{tabular}{rrrrrr}
\toprule
Row & $y$ & $x_1$ & $x_2$ & $x_3$ & $x_4$\\
\midrule
1 & -2.191988 & -1.479972 & -1.698022 & -0.465721 & -0.818119\\
2 & -0.502206 & -0.662950 &  0.653037 &  0.738444 &  0.336596\\
3 &  0.302931 & -0.386687 & -0.993984 & -1.748198 & -1.520414\\
\bottomrule
\end{tabular}
\end{center}
Let $X=(\ones,x)$, with $n=200$ rows and $k=5$ columns, and let
$\widehat r=y-X\widehat\beta$. The model-based variance is
$\widehat\sigma^2(X\trans X)^{-1}$ with
$\widehat\sigma^2=\widehat r\trans\widehat r/(n-k)$, and the robust (HC0)
variance is $(X\trans X)^{-1}X\trans\widehat DX(X\trans X)^{-1}$ with
$\widehat D=\operatorname{diag}(\widehat r_1^2,\ldots,\widehat r_n^2)$.
\begin{center}
\begin{tabular}{lrrr}
\toprule
Coefficient & Estimate & Model SE & HC0 SE\\
\midrule
Intercept & 0.954681 & 0.098797 & 0.097826\\
$x_1$     & 0.683716 & 0.163978 & 0.142299\\
$x_2$     & 0.739820 & 0.209943 & 0.225490\\
$x_3$     & 0.256060 & 0.220531 & 0.217935\\
$x_4$     & 0.868421 & 0.183722 & 0.235053\\
\bottomrule
\end{tabular}
\end{center}
 
\textbf{(b)} For each coefficient $j$, residualize the column $X_j$ on the
other four columns to get $R_j$. By Exercise 1.7 (which applies to any column
after reordering),
\[
\widehat\beta_j=\frac{R_j\trans y}{R_j\trans R_j},
\qquad
\mathrm{SE}_{\mathrm{mod},j}=\sqrt{\frac{\widehat\sigma^2}{R_j\trans R_j}},
\qquad
\mathrm{SE}_{\mathrm{HC0},j}=\sqrt{\frac{R_j\trans\widehat DR_j}{(R_j\trans R_j)^2}}.
\]
Here $\widehat\sigma^2$ uses the full-model degrees of freedom $n-k=195$; by
Exercise 1.7(a), using $n-1$ from the residual-on-residual fit would instead
shrink each model SE by a factor of $\sqrt{195/199}$. These formulas reproduce
the table above, with a maximum absolute difference of $6.44\times10^{-15}$.
 
\textbf{(c)} Let $a=(0,1,1,1)\trans$ and $S=a\trans X=X_2+X_3+X_4$, so
$Y=X_1+S_+^2/3+\varepsilon$ with $\varepsilon$ independent of $X$. The target
minimizes $\E(Y-\beta_0-X\trans b)^2$. Since $\E X=0$, its normal equations
give $\beta_0^*=\E Y$ and $b^*=\Sigma^{-1}\operatorname{Cov}(X,Y)$.
 
\emph{Intercept.} $S\sim N(0,v_S)$ with
\[
v_S=3+2\{\operatorname{Cov}(X_2,X_3)+\operatorname{Cov}(X_2,X_4)
+\operatorname{Cov}(X_3,X_4)\}=3+2(0.8+0.64+0.8)=7.48 .
\]
By symmetry of $S$, $\E S_+^2=\frac12\E S^2=v_S/2$, so
$\beta_0^*=\E(S_+^2)/3=v_S/6=1.2466667$.
 
\emph{Slopes.} $\operatorname{Cov}(X,Y)=\Sigma e_1+\frac13\E(XS_+^2)$. Because
$(X,S)$ is jointly normal, $\E(X\mid S)=(\Sigma a/v_S)\,S$, so
\[
\E(XS_+^2)=\frac{\Sigma a}{v_S}\,\E\bigl(S^3\ind\{S>0\}\bigr)
=\frac{\Sigma a}{v_S}\,v_S^{3/2}\sqrt{2/\pi},
\]
using $\E(Z^3\ind\{Z>0\})=\sqrt{2/\pi}$ for $Z\sim N(0,1)$. Therefore
\[
b^*=e_1+\lambda a,\qquad
\lambda=\frac{\sqrt{v_S}}{3}\sqrt{\frac2\pi}=0.7273938,
\]
and
$\beta^*=(1.2466667,\ 1,\ 0.7273938,\ 0.7273938,\ 0.7273938)\trans$.
 
The following R functions implement the direct and residual-on-residual
calculations used above.
\begin{lstlisting}
ols_stats <- function(x, y) {
  X <- cbind(1, x)
  n <- nrow(X)
  k <- ncol(X)
  XtX_inv <- solve(crossprod(X))
  beta_hat <- drop(XtX_inv %*% crossprod(X, y))
  residual <- drop(y - X %*% beta_hat)
 
  V_model <- sum(residual^2) / (n - k) * XtX_inv
  V_robust <- XtX_inv %*%
    crossprod(X, X * residual^2) %*% XtX_inv
 
  list(beta = beta_hat, residual = residual,
       V_model = V_model, V_robust = V_robust, X = X)
}
 
ror_stats <- function(X, y) {
  n <- nrow(X)
  k <- ncol(X)
  full <- ols_stats(X[, -1, drop = FALSE], y)
  out <- matrix(NA_real_, k, 3)
 
  for (j in seq_len(k)) {
    Z <- X[, j]
    W <- X[, -j, drop = FALSE]
    RZ <- drop(Z - W %*%
      solve(crossprod(W), crossprod(W, Z)))
 
    out[j, 1] <- sum(RZ * y) / sum(RZ^2)
    out[j, 2] <- sqrt(
      sum(full$residual^2) / (n - k) / sum(RZ^2))
    out[j, 3] <- sqrt(
      sum(RZ^2 * full$residual^2) / sum(RZ^2)^2)
  }
  colnames(out) <- c("estimate", "model_se", "robust_se")
  out
}
 
fit <- ols_stats(x, y)
standard <- cbind(
  estimate = fit$beta,
  model_se = sqrt(diag(fit$V_model)),
  robust_se = sqrt(diag(fit$V_robust)))
ror <- ror_stats(fit$X, y)
standard
ror
max(abs(standard - ror))
\end{lstlisting}
\end{solution}
 
%==============================================================================
\question \textbf{Exercise 1.9.}
Using the random-design setting of Exercise 1.8, run 2000 repeated
simulations and report:
 
\begin{parts}
\part Monte Carlo biases and standard deviations of the five estimated
coefficients.
 
\part Square roots of the Monte Carlo means of model-based variances and the
$90\%$ and $95\%$ Wald coverage proportions.
 
\part Square roots of the Monte Carlo means of model-robust variances and the
$90\%$ and $95\%$ Wald coverage proportions.
 
\part Summarize the conclusions from the simulation.
\end{parts}
 
\begin{solution}
The target is $\beta^*=(1.2466667,\ 1,\ 0.7273938,\ 0.7273938,\ 0.7273938)\trans$
from Exercise 1.8. There are 2000 replications after \texttt{set.seed(0)}; each draws new $X$ and $Y$, and
the $90\%$ and $95\%$ Wald intervals use the normal critical values $1.645$ and
$1.960$. ``RMV'' is the square root of the Monte Carlo mean of the estimated
variances.
 
\textbf{(a)--(b)} Model-based results (2000 replications):
\begin{center}
\small
\setlength{\tabcolsep}{4pt}
\begin{tabular}{lrrrrrr}
\toprule
Coefficient & Target & Bias & MC SD & Model RMV & Cov.\ 90\% & Cov.\ 95\%\\
\midrule
Intercept & 1.246667 & -0.016471 & 0.144232 & 0.141715 & 0.8655 & 0.9310\\
$x_1$     & 1.000000 & -0.001998 & 0.229680 & 0.237606 & 0.9130 & 0.9585\\
$x_2$     & 0.727394 & -0.014301 & 0.315431 & 0.304206 & 0.8755 & 0.9345\\
$x_3$     & 0.727394 &  0.002095 & 0.311521 & 0.303437 & 0.8885 & 0.9425\\
$x_4$     & 0.727394 & -0.003744 & 0.253048 & 0.236669 & 0.8750 & 0.9280\\
\bottomrule
\end{tabular}
\end{center}
 
\textbf{(c)} HC0 results (same 2000 replications):
\begin{center}
\begin{tabular}{lrrr}
\toprule
Coefficient & HC0 RMV & Cov.\ 90\% & Cov.\ 95\%\\
\midrule
Intercept & 0.140476 & 0.8610 & 0.9280\\
$x_1$     & 0.229364 & 0.9025 & 0.9545\\
$x_2$     & 0.307878 & 0.8795 & 0.9315\\
$x_3$     & 0.309789 & 0.8830 & 0.9385\\
$x_4$     & 0.250629 & 0.8805 & 0.9295\\
\bottomrule
\end{tabular}
\end{center}
 
\textbf{(d)} The linear model is misspecified, since the true mean contains
$S_+^2/3$, but OLS still estimates the projection target $\beta^*$.
 
\emph{Bias.} The biases are small compared with the MC SDs. The intercept bias
($-0.0165$) is about five Monte Carlo standard errors
($0.144/\sqrt{2000}\approx0.0032$) from zero, which is strong Monte Carlo evidence
of finite-sample bias in this data-generating process: here OLS is consistent
for $\beta^*$ but not exactly unbiased at $n=200$. The slope biases are within about two Monte Carlo
standard errors of zero.
 
\emph{Standard errors.} The error around the projection,
$Y-X\trans\beta^*=\varepsilon+\{m(X)-X\trans\beta^*\}$, combines the noise
$\varepsilon$ (variance $0.25$) with the approximation error of the nonlinear
mean, and the second part makes its variance depend on $X$. For $x_1$ and $x_4$
the HC0 RMVs ($0.229$, $0.251$) match the MC SDs ($0.230$, $0.253$) closely,
while the model-based RMVs are too large for $x_1$ ($0.238$) and too small for
$x_4$ ($0.237$). This pattern is consistent with heteroskedasticity; the
direction of the model-based error can differ across coefficients. For $x_2$
and $x_3$, both estimators understate the MC SD by up to about $4\%$, with HC0
closer.
 
\emph{Coverage.} A coverage proportion $\pi$ has Monte Carlo standard error
$\sqrt{\pi(1-\pi)/2000}$, which is about $0.005$ at $\pi=0.95$ and $0.007$ at
$\pi=0.90$. So coverage near $0.93$ is genuinely below nominal. Both
methods undercover slightly for the intercept, $x_2$, $x_3$, and $x_4$: HC0
has no small-sample correction, and the intercept is also biased. For $x_1$
the model-based intervals overcover because its SE is too large, while the
HC0 intervals are close to nominal.
 
The R code for the simulation is
\begin{lstlisting}
summarize_mc <- function(B, V_model, V_robust, target) {
  z90 <- qnorm(.95)
  z95 <- qnorm(.975)
  coverage <- function(V, z) {
    colMeans(abs(sweep(B, 2, target, "-")) <= z * sqrt(V))
  }
 
  data.frame(
    target = target,
    mean = colMeans(B),
    bias = colMeans(B) - target,
    mc_sd = apply(B, 2, sd),
    model_root_mean_var = sqrt(colMeans(V_model)),
    model_cov90 = coverage(V_model, z90),
    model_cov95 = coverage(V_model, z95),
    robust_root_mean_var = sqrt(colMeans(V_robust)),
    robust_cov90 = coverage(V_robust, z90),
    robust_cov95 = coverage(V_robust, z95))
}
 
set.seed(0)
m <- 2000
coef_names <- c("(Intercept)", paste0("x", 1:4))
a <- c(0, 1, 1, 1)
vS <- drop(t(a) %*% msig %*% a)
lambda <- sqrt(vS) * sqrt(2 / pi) / 3
target_random <- c(vS / 6, 1, rep(lambda, 3))
 
B <- V_model <- V_robust <- matrix(NA_real_, m, 5)
for (s in seq_len(m)) {
  x_s <- mvrnorm(n, mu, msig)
  y_s <- x_s[, 1] +
    pmax(0, rowSums(x_s[, 2:4, drop = FALSE]))^2 / 3 +
    rnorm(n, 0, sig)
 
  fit_s <- ols_stats(x_s, y_s)
  B[s, ] <- fit_s$beta
  V_model[s, ] <- diag(fit_s$V_model)
  V_robust[s, ] <- diag(fit_s$V_robust)
}
 
random_results <- summarize_mc(
  B, V_model, V_robust, target_random)
rownames(random_results) <- coef_names
random_results
\end{lstlisting}
\end{solution}
 
%==============================================================================
\question \textbf{Exercise 1.10.}
Use the setting of Exercise 1.8, but hold the generated $x$-matrix fixed.
 
\begin{parts}
\part Find analytical expressions and numerical values for all fixed-design
target coefficients.
 
\part Using 2000 repeated simulations with $x$ fixed, report the Monte Carlo
biases and standard deviations of the five estimated coefficients.
 
\part Report the square roots of the Monte Carlo means of model-based
variances and the $90\%$ and $95\%$ coverage proportions.
 
\part Report the corresponding model-robust quantities and summarize the
simulation results.
\end{parts}
 
\begin{solution}
\textbf{(a)} Fix the $x$ generated right after \texttt{set.seed(0)} (the same
$x$ as in Exercise 1.8) and let $X=(\ones,x)$. Given $X$, we have
$Y=m+\varepsilon$ with
\[
m_i=x_{i1}+\frac{\{\max(0,x_{i2}+x_{i3}+x_{i4})\}^2}{3},
\qquad
\varepsilon\sim N(0,\sigma^2I_n),\quad \sigma=0.5 .
\]
So
\[
\widehat\beta=(X\trans X)^{-1}X\trans m+(X\trans X)^{-1}X\trans\varepsilon,
\]
and the target is
$\beta_X^*=\E(\widehat\beta\mid X)=(X\trans X)^{-1}X\trans m$, the projection
of the 200 true means onto this particular design:
\begin{center}
\begin{tabular}{lr}
\toprule
Coefficient & Fixed-design target\\
\midrule
Intercept & 0.983946\\
$x_1$     & 0.650521\\
$x_2$     & 0.815593\\
$x_3$     & 0.103363\\
$x_4$     & 0.973887\\
\bottomrule
\end{tabular}
\end{center}
These differ from the random-design targets because they depend on the
realized $x$.
 
\textbf{(b)--(c)} Model-based results (2000 replications after
\texttt{set.seed(0)} and the draw of $x$). Only $\varepsilon$ is redrawn in each
replication, and coverage is measured against $\beta_X^*$.
\begin{center}
\small
\setlength{\tabcolsep}{4pt}
\begin{tabular}{lrrrrrr}
\toprule
Coefficient & Target & Bias & MC SD & Model RMV & Cov.\ 90\% & Cov.\ 95\%\\
\midrule
Intercept & 0.983946 &  0.000224 & 0.035434 & 0.097466 & 1.0000 & 1.0000\\
$x_1$     & 0.650521 & -0.001225 & 0.059715 & 0.161769 & 1.0000 & 1.0000\\
$x_2$     & 0.815593 &  0.000341 & 0.075042 & 0.207115 & 1.0000 & 1.0000\\
$x_3$     & 0.103363 &  0.000391 & 0.078367 & 0.217560 & 1.0000 & 1.0000\\
$x_4$     & 0.973887 &  0.000558 & 0.067411 & 0.181247 & 1.0000 & 1.0000\\
\bottomrule
\end{tabular}
\end{center}
 
\textbf{(d)} HC0 results (same replications):
\begin{center}
\begin{tabular}{lrrr}
\toprule
Coefficient & HC0 RMV & Cov.\ 90\% & Cov.\ 95\%\\
\midrule
Intercept & 0.096745 & 1.0000 & 1.0000\\
$x_1$     & 0.139671 & 0.9995 & 1.0000\\
$x_2$     & 0.229211 & 1.0000 & 1.0000\\
$x_3$     & 0.225230 & 1.0000 & 1.0000\\
$x_4$     & 0.238594 & 1.0000 & 1.0000\\
\bottomrule
\end{tabular}
\end{center}
Given $X$, OLS is exactly unbiased for $\beta_X^*$, and by the decomposition
in (a),
\[
\operatorname{Var}(\widehat\beta\mid X)
=(X\trans X)^{-1}X\trans(\sigma^2I)X(X\trans X)^{-1}
=\sigma^2(X\trans X)^{-1}=0.25\,(X\trans X)^{-1},
\]
which is what the MC SDs estimate. Three quantities should be kept apart
here: the noise $\varepsilon$, with variance $0.25$; the approximation error
$m-X\beta_X^*$ of the best linear fit, which is fixed once $X$ is fixed;
and the OLS residual $\widehat r$, which contains both. The variance
estimators are built from $\widehat r$. With $H=X(X\trans X)^{-1}X\trans$,
\[
\widehat r=(I-H)m+(I-H)\varepsilon .
\]
Because the mean is nonlinear, $m\notin\col(X)$, so the first term,
$(I-H)m=m-X\beta_X^*$, is a fixed, nonzero approximation error. It inflates both estimators:
$\E(\widehat\sigma^2\mid X)=\sigma^2+\normtwo{(I-H)m}^2/(n-k)$, and
$\E(\widehat r_i^2\mid X)=\{((I-H)m)_i\}^2+\sigma^2(1-h_{ii})$. Both treat the
fixed approximation error as if it were noise, so the model RMV is approximately $2.7$
times the MC SD for each coefficient and coverage is essentially $100\%$.
 
This does not contradict Exercise 1.9. Under random design the approximation
error changes from sample to sample and really is part of the variability of
$\widehat\beta$; holding one design fixed removes it.
 
The R code is
\begin{lstlisting}
set.seed(0)
x_fixed <- mvrnorm(n, mu, msig)
X_fixed <- cbind(1, x_fixed)
mean_y_fixed <- x_fixed[, 1] +
  pmax(0, rowSums(x_fixed[, 2:4, drop = FALSE]))^2 / 3
 
target_fixed <- drop(
  solve(crossprod(X_fixed),
        crossprod(X_fixed, mean_y_fixed)))
 
m <- 2000
B <- V_model <- V_robust <- matrix(NA_real_, m, 5)
for (s in seq_len(m)) {
  y_s <- mean_y_fixed + rnorm(n, 0, sig)
  fit_s <- ols_stats(x_fixed, y_s)
 
  B[s, ] <- fit_s$beta
  V_model[s, ] <- diag(fit_s$V_model)
  V_robust[s, ] <- diag(fit_s$V_robust)
}
 
fixed_results <- summarize_mc(
  B, V_model, V_robust, target_fixed)
rownames(fixed_results) <- coef_names
 
target_fixed
fixed_results
\end{lstlisting}
\end{solution}
 
\end{questions}
\end{document}
 
